\documentclass[11pt,a4paper]{article}

\usepackage[utf8]{inputenc}
\usepackage{amsmath,amssymb,amsfonts}
\usepackage{booktabs,tabularx,longtable}
\usepackage{enumitem}
\usepackage{geometry}
\usepackage{hyperref}
\usepackage{xcolor}
\usepackage{listings}
\usepackage{fancyhdr}
\usepackage{titlesec}

\geometry{margin=2.3cm,top=2.8cm,bottom=2.8cm}

\hypersetup{
  colorlinks=true,
  linkcolor=blue!60!black,
  citecolor=blue!60!black,
  urlcolor=blue!60!black,
  pdftitle={ISLS v2.1 --- Multi-Pass Rendering},
  pdfauthor={Sebastian Klemm},
}

\lstset{
  basicstyle=\ttfamily\footnotesize,
  breaklines=true,
  frame=single,
  backgroundcolor=\color{gray!8},
  keywordstyle=\color{blue!70!black}\bfseries,
  stringstyle=\color{red!60!black},
  commentstyle=\color{green!50!black}\itshape,
  numbers=left,
  numberstyle=\tiny\color{gray},
  numbersep=6pt,
  tabsize=4,
  showstringspaces=false,
}

\titleformat{\section}{\Large\bfseries}{\thesection}{1em}{}
\titleformat{\subsection}{\large\bfseries}{\thesubsection}{1em}{}
\titleformat{\subsubsection}{\normalsize\bfseries}{\thesubsubsection}{1em}{}

\pagestyle{fancy}
\fancyhf{}
\fancyhead[L]{\small ISLS v2.1 --- Multi-Pass Rendering}
\fancyhead[R]{\small Klemm, 2026}
\fancyfoot[C]{\thepage}

\begin{document}

\begin{titlepage}
\begin{center}
\vspace*{1.5cm}
{\Huge\bfseries ISLS v2.1}\\[0.4cm]
{\LARGE\bfseries Multi-Pass Rendering}\\[0.3cm]
{\LARGE\bfseries with Universal Crystals}\\[0.3cm]
{\LARGE\bfseries and LLM-Driven Domain Intelligence}\\[1.2cm]
{\Large Specification v1.0}\\[1.5cm]
{\large Sebastian Klemm}\\[0.3cm]
{\small March 2026}\\[2cm]

\begin{tabular}{ll}
\textbf{Repository:} & \texttt{graphity/isls} \\
\textbf{Builds on:} & ISLS v2.0 (Hypercube Decomposer) \\
\textbf{New crate:} & \texttt{isls-renderloop} \\
\textbf{Modified:} & \texttt{isls-decomposer}, \texttt{isls-blueprint},
  \texttt{isls-orchestrator}, \texttt{isls-cli} \\
\textbf{Requires:} & OpenAI API key (for LLM passes) \\
\end{tabular}

\vfill
{\small
v2.0 generates in one pass: templates produce structure, done.\\
v2.1 renders in multiple passes: structure, then logic, then edge cases,\\
then polish --- each pass deeper, each pass cheaper, converging.\\[0.3cm]
The crystal system is restructured: domain-specific patterns are\\
replaced by \textbf{universal architectural patterns} that apply to\\
ANY application. The system does not memorize warehouses.\\
It learns what software \emph{is}.\\[0.3cm]
The LLM is not a fallback. It is the creative engine that translates\\
human domains into architectural structure.
}
\end{center}
\end{titlepage}

\tableofcontents
\newpage

% ============================================================
\section{Why v2.0 Is Not Enough}

v2.0 produces 60 files and 5,120 LOC in 0.8 seconds with zero Oracle
calls. Impressive. But the generated code is template code. Every
function body follows the same pattern. Every service does the same
thing: validate, query, return. The business logic is structurally
correct but semantically empty.

And worse: the blueprint system learns \textbf{domain-specific}
patterns. After generating 5 warehouse systems, it is excellent at
warehouses and useless at everything else. The crystals are
``Warehouse Product with 17 fields'' and ``Warehouse Order
fulfillment.'' A music streaming service cannot use these. The system
has memorized, not learned.

v2.1 fixes both problems:

\begin{enumerate}
\item \textbf{Multi-pass rendering}: Structure first (offline),
  then domain logic (LLM), then edge cases (LLM), then polish
  (LLM). Each pass deeper. Each pass cheaper. Converging.
\item \textbf{Universal crystals}: The system does not memorize
  ``warehouse product.'' It learns ``entity with inventory tracking
  needs: quantity field, reorder threshold, stock movement logging.''
  This applies to warehouse products, pharmacy inventory, library
  books, and spare parts. One crystal, infinite domains.
\end{enumerate}

% ============================================================
\section{Universal Crystal Architecture}

\subsection{Crystal Levels}

Every crystal has a level indicating its abstraction:

\begin{lstlisting}[language=C]
/// The abstraction level of a crystal.
/// Higher levels are more universal, more reusable, more valuable.
pub enum CrystalLevel {
    /// Level 0: Universal architectural patterns.
    /// Apply to ANY software, ANY domain.
    /// Examples: BelongsTo, HasMany, CRUD, Pagination, Auth,
    ///           StateMachine, EventEmitter, Pipeline.
    /// These NEVER become stale.
    Universal,

    /// Level 1: Architectural pattern families.
    /// Apply to classes of applications.
    /// Examples: "REST API with auth", "Event-driven pipeline",
    ///           "CRUD with soft delete", "Multi-tenant SaaS".
    Architectural,

    /// Level 2: Structural patterns.
    /// Apply across many domains but are more specific.
    /// Examples: "Entity with inventory tracking",
    ///           "Entity with state machine and transitions",
    ///           "Service with external API integration".
    Structural,

    /// Level 3: Domain-specific patterns.
    /// Apply only within one domain.
    /// Examples: "Warehouse stock adjustment",
    ///           "E-commerce checkout flow".
    /// These are LEAST valuable for reuse and should be
    /// the LAST to be consulted.
    DomainSpecific,
}
\end{lstlisting}

\subsection{Crystal Structure (Revised)}

\begin{lstlisting}[language=C]
pub struct Crystal {
    pub id: String,                 // SHA-256
    pub level: CrystalLevel,
    pub pattern_name: String,       // "belongs_to", "state_machine", etc.
    pub description: String,
    pub confidence: f64,

    /// What this crystal knows about structure
    pub structure: StructuralKnowledge,

    /// What this crystal knows about implementation
    pub implementation: ImplementationKnowledge,

    /// Usage statistics
    pub stats: CrystalStats,

    /// Evidence chain
    pub evidence: Vec<String>,
}

/// Structural knowledge: WHAT should be built.
/// This is the template-level information.
pub struct StructuralKnowledge {
    /// What types/fields/signatures are needed
    pub components: Vec<ComponentSpec>,
    /// How components relate to each other
    pub relationships: Vec<RelSpec>,
    /// Constraints that must be satisfied
    pub constraints: Vec<String>,
}

pub struct ComponentSpec {
    pub kind: String,       // "struct", "function", "endpoint", "migration"
    pub name_pattern: String,  // "{entity}_service", "list_{entities}"
    pub parameters: Vec<ParamSpec>,
    pub returns: Option<String>,
    pub description: String,
}

/// Implementation knowledge: HOW it should be built.
/// This is what the LLM teaches us after generation.
pub struct ImplementationKnowledge {
    /// Common patterns in the function body
    pub body_patterns: Vec<String>,
    /// Common error cases to handle
    pub error_cases: Vec<String>,
    /// Common edge cases discovered
    pub edge_cases: Vec<String>,
    /// Performance considerations
    pub perf_notes: Vec<String>,
    /// Security considerations
    pub security_notes: Vec<String>,
    /// Test scenarios that should exist
    pub test_scenarios: Vec<String>,
}

pub struct CrystalStats {
    pub times_used: u64,
    pub times_succeeded: u64,
    pub domains_used_in: Vec<String>,  // ["warehouse", "ecommerce", "hospital"]
    pub last_used: String,
    pub avg_tokens_saved: f64,  // estimated tokens NOT spent because of this crystal
}
\end{lstlisting}

\subsection{Built-in Universal Crystals (Level 0)}

The system ships with these pre-defined crystals. They are never
deleted and cannot fall below confidence 0.9.

\begin{center}
\small
\begin{tabular}{lp{8cm}}
\toprule
\textbf{Crystal} & \textbf{What It Encodes} \\
\midrule
\texttt{belongs\_to} & Entity A references Entity B via foreign key. Needs: FK field,
  FK index, ON DELETE policy, join query, nested serialization option. \\
\texttt{has\_many} & Entity A owns many Entity B. Needs: reverse FK in B,
  list-children query, cascade delete option, count method. \\
\texttt{crud\_entity} & Entity with Create/Read/Update/Delete. Needs: model struct,
  CreatePayload, UpdatePayload, validation, 5 query functions, 5 API endpoints. \\
\texttt{state\_machine} & Entity with discrete states and valid transitions. Needs:
  state enum, transition validation function, status history logging. \\
\texttt{pagination} & List endpoint with page/per\_page/sort/search. Needs:
  PaginationParams, PaginatedResponse, LIMIT/OFFSET query, count query. \\
\texttt{jwt\_auth} & JWT authentication flow. Needs: Claims struct, create/verify
  functions, AuthUser extractor, role hierarchy, login/register endpoints. \\
\texttt{error\_handling} & Typed error system. Needs: AppError enum, Display,
  ResponseError, From<sqlx::Error>, proper HTTP status codes. \\
\texttt{soft\_delete} & Entity is never physically deleted, only marked inactive.
  Needs: is\_active/deleted\_at field, filtered queries, restore function. \\
\texttt{audit\_trail} & Track who changed what when. Needs: audit table, trigger
  or service-level logging, query audit history. \\
\texttt{event\_emitter} & Action produces domain events. Needs: Event enum,
  EventBus, subscriber registration, async dispatch. \\
\texttt{pipeline} & Data flows through stages: input $\to$ validate $\to$ transform
  $\to$ persist $\to$ respond. Needs: stage functions, error at each stage. \\
\texttt{config\_from\_env} & Configuration loaded from environment variables with
  defaults. Needs: Config struct, env parsing, validation. \\
\texttt{health\_check} & System health endpoint. Needs: /health route, DB ping,
  uptime, version. \\
\texttt{rate\_limiting} & API rate limiting per user/IP. Needs: middleware,
  sliding window counter, 429 response. \\
\texttt{background\_job} & Async task processed outside request cycle. Needs:
  job queue, worker loop, retry logic, failure handling. \\
\bottomrule
\end{tabular}
\end{center}

These 15 crystals cover the structural skeleton of virtually any
web application. They are the ``laws of physics'' of the ISLS
universe. Domain-specific behavior lives in the LLM passes, not
in these crystals.

\subsection{Crystal Matching (Revised)}

\begin{lstlisting}[language=C]
impl BlueprintRegistry {
    /// Find matching crystals for a dimension.
    /// Search order: Universal -> Architectural -> Structural -> Domain.
    /// Returns ALL matches above threshold, sorted by level (most universal first).
    pub fn find_matches(
        &self,
        dimension: &Dimension,
        min_confidence: f64,
    ) -> Vec<&Crystal> {
        let mut matches = Vec::new();

        // Level 0: Universal — always check
        for c in self.by_level(CrystalLevel::Universal) {
            if c.matches_dimension(dimension) && c.confidence >= min_confidence {
                matches.push(c);
            }
        }

        // Level 1: Architectural
        for c in self.by_level(CrystalLevel::Architectural) {
            if c.matches_dimension(dimension) && c.confidence >= min_confidence {
                matches.push(c);
            }
        }

        // Level 2: Structural (only if no Universal/Architectural match)
        if matches.is_empty() {
            for c in self.by_level(CrystalLevel::Structural) {
                if c.matches_dimension(dimension) && c.confidence >= min_confidence {
                    matches.push(c);
                }
            }
        }

        // Level 3: Domain-specific (only as last resort)
        // These are least likely to generalize
        if matches.is_empty() {
            for c in self.by_level(CrystalLevel::DomainSpecific) {
                if c.matches_dimension(dimension) && c.confidence >= min_confidence {
                    matches.push(c);
                }
            }
        }

        matches
    }
}
\end{lstlisting}

\subsection{How Crystals Gain Implementation Knowledge}

When the LLM generates code in a render pass, and that code survives
verification, the system extracts implementation knowledge and adds
it to the matching crystal:

\begin{lstlisting}[language=C]
/// After a successful LLM pass, learn from what the LLM produced.
fn learn_from_llm_output(
    crystal: &mut Crystal,
    generated_code: &str,
    pass_depth: u32,
) {
    match pass_depth {
        1 => {
            // Logic pass: learn body patterns
            let patterns = extract_code_patterns(generated_code);
            for p in patterns {
                if !crystal.implementation.body_patterns.contains(&p) {
                    crystal.implementation.body_patterns.push(p);
                }
            }
        }
        2 => {
            // Edge case pass: learn error cases
            let errors = extract_error_handling(generated_code);
            for e in errors {
                if !crystal.implementation.error_cases.contains(&e) {
                    crystal.implementation.error_cases.push(e);
                }
            }
            let edges = extract_edge_cases(generated_code);
            for e in edges {
                if !crystal.implementation.edge_cases.contains(&e) {
                    crystal.implementation.edge_cases.push(e);
                }
            }
        }
        3 => {
            // Polish pass: learn test scenarios, security notes
            let tests = extract_test_patterns(generated_code);
            for t in tests {
                if !crystal.implementation.test_scenarios.contains(&t) {
                    crystal.implementation.test_scenarios.push(t);
                }
            }
        }
        _ => {}
    }
}
\end{lstlisting}

Over time, the universal crystals accumulate implementation knowledge
from every domain they participate in. The \texttt{state\_machine}
crystal learns error cases from warehouse orders, medical patient
workflows, and project task boards. Each domain teaches it something
new. The crystal becomes wiser without becoming domain-specific.

% ============================================================
\section{Multi-Pass Render Loop}

\subsection{Pass Definitions}

\begin{lstlisting}[language=C]
pub struct RenderPass {
    /// Pass depth: 0=structure, 1=logic, 2=edge cases, ...
    pub depth: u32,
    /// What this pass does
    pub pass_type: PassType,
    /// Which files/functions this pass operates on
    pub scope: PassScope,
    /// Maximum tokens budget for this pass
    pub token_budget: u64,
    /// Stop when change rate drops below this threshold
    pub convergence_threshold: f64,
}

pub enum PassType {
    /// Pass 0: Structure generation from templates + universal crystals.
    /// Offline. Zero tokens. Produces: files, structs, function signatures,
    /// SQL schema, API routes, frontend pages.
    Structure,

    /// Pass 1: Domain logic injection.
    /// LLM fills function bodies with actual business logic.
    /// Guided by: crystal's StructuralKnowledge + TOML descriptions.
    DomainLogic,

    /// Pass 2: Edge case hardening.
    /// LLM reviews generated code and adds error handling,
    /// boundary checks, race condition guards.
    /// Guided by: crystal's ImplementationKnowledge.error_cases.
    EdgeCases,

    /// Pass 3: Integration pass.
    /// LLM checks cross-module consistency: do API responses
    /// match frontend expectations? Do queries match models?
    Integration,

    /// Pass 4: Test generation.
    /// LLM writes tests based on crystal's test_scenarios
    /// plus domain-specific behavior from earlier passes.
    TestGeneration,

    /// Pass 5: Polish.
    /// LLM improves: error messages, logging, documentation,
    /// code comments, naming conventions.
    Polish,
}

pub enum PassScope {
    /// All generated files
    All,
    /// Only files in a specific layer
    Layer(String),  // "services", "api", "frontend"
    /// Only specific files
    Files(Vec<String>),
    /// Only functions matching a pattern
    Functions(String),  // "adjust_stock", "*_service"
}
\end{lstlisting}

\subsection{The Render Loop}

\begin{lstlisting}[language=C]
/// The multi-pass render loop.
/// Takes v2.0 decomposer output (structure) and renders it
/// through multiple LLM passes until convergence.
pub struct RenderLoop {
    /// Available LLM oracle
    oracle: Box<dyn Oracle>,
    /// Crystal registry
    crystals: BlueprintRegistry,
    /// Barbara reader for verification
    reader: Reader,
    /// Barbara topology analyzer
    topo: CodeTopoAnalyzer,
    /// Render pass configuration
    passes: Vec<RenderPass>,
    /// Statistics
    stats: RenderStats,
}

pub struct RenderStats {
    pub passes_executed: u32,
    pub total_tokens_used: u64,
    pub tokens_per_pass: Vec<u64>,
    pub files_modified_per_pass: Vec<usize>,
    pub convergence_per_pass: Vec<f64>,
    pub crystals_updated: usize,
}

impl RenderLoop {
    /// Default pass sequence for a full-stack application.
    pub fn default_passes() -> Vec<RenderPass> {
        vec![
            RenderPass {
                depth: 0,
                pass_type: PassType::Structure,
                scope: PassScope::All,
                token_budget: 0,  // offline
                convergence_threshold: 0.0,
            },
            RenderPass {
                depth: 1,
                pass_type: PassType::DomainLogic,
                scope: PassScope::Layer("services".into()),
                token_budget: 50_000,
                convergence_threshold: 0.05,
            },
            RenderPass {
                depth: 2,
                pass_type: PassType::EdgeCases,
                scope: PassScope::Layer("services".into()),
                token_budget: 20_000,
                convergence_threshold: 0.03,
            },
            RenderPass {
                depth: 3,
                pass_type: PassType::Integration,
                scope: PassScope::All,
                token_budget: 15_000,
                convergence_threshold: 0.02,
            },
            RenderPass {
                depth: 4,
                pass_type: PassType::TestGeneration,
                scope: PassScope::Layer("tests".into()),
                token_budget: 20_000,
                convergence_threshold: 0.05,
            },
            RenderPass {
                depth: 5,
                pass_type: PassType::Polish,
                scope: PassScope::All,
                token_budget: 10_000,
                convergence_threshold: 0.01,
            },
        ]
    }

    /// Execute the full render loop.
    pub fn render(
        &mut self,
        mut artifacts: Vec<Artifact>,  // from v2.0 decomposer (Pass 0)
        app_spec: &AppSpec,
    ) -> Result<Vec<Artifact>> {
        self.stats.passes_executed = 1;  // Pass 0 already done

        for pass in &self.passes[1..] {  // skip Pass 0 (structure)
            let before = snapshot_artifacts(&artifacts);
            let tokens_before = self.stats.total_tokens_used;

            // Execute pass
            artifacts = self.execute_pass(pass, artifacts, app_spec)?;

            // Measure convergence
            let change_rate = measure_change_rate(&before, &artifacts);
            let tokens_used = self.stats.total_tokens_used - tokens_before;

            self.stats.passes_executed += 1;
            self.stats.tokens_per_pass.push(tokens_used);
            self.stats.convergence_per_pass.push(change_rate);
            self.stats.files_modified_per_pass.push(
                count_modified(&before, &artifacts)
            );

            // Check convergence
            if change_rate < pass.convergence_threshold {
                // This pass converged. Move to next.
                continue;
            }

            // Check token budget
            if tokens_used > pass.token_budget {
                // Budget exhausted for this pass. Move to next.
                continue;
            }
        }

        // After all passes: crystallize what we learned
        self.crystallize_learnings(&artifacts);

        Ok(artifacts)
    }
}
\end{lstlisting}

\subsection{Pass 0: Structure (Offline)}

This is exactly what v2.0 does today. The Hypercube Decomposer runs,
templates produce code. Zero tokens. The output is:

\begin{itemize}
\item All files created with correct names, paths, module structure
\item All structs with all fields (from domain knowledge)
\item All function signatures (correct parameters, return types)
\item SQL migrations with complete schema
\item API routes with correct HTTP methods and paths
\item Frontend pages with layout structure
\item Config, Dockerfile, docker-compose, README
\end{itemize}

What is \emph{missing} after Pass 0:
\begin{itemize}
\item Function bodies are template boilerplate (CRUD only)
\item Business rules are comments, not code
\item Error handling is minimal
\item No edge case handling
\item Tests are structural stubs
\item Frontend has no actual interaction logic
\end{itemize}

\subsection{Pass 1: Domain Logic (LLM)}

For each service function that has business rule comments or
template boilerplate bodies, construct a prompt:

\begin{lstlisting}[language=C]
fn build_domain_logic_prompt(
    function: &FunctionArtifact,
    crystal: &Crystal,
    app_spec: &AppSpec,
) -> String {
    format!(
        "You are filling in the business logic for a Rust function.\n\
         \n\
         ## Function Signature (DO NOT CHANGE)\n\
         ```rust\n{}\n```\n\
         \n\
         ## Available Types\n\
         ```rust\n{}\n```\n\
         \n\
         ## Available Database Queries\n\
         {}\n\
         \n\
         ## Business Rule to Implement\n\
         {}\n\
         \n\
         ## Known Patterns from Similar Functions\n\
         {}\n\
         \n\
         ## Known Error Cases to Handle\n\
         {}\n\
         \n\
         Return ONLY the function body (the code between {{ and }}).\n\
         Use the exact types and queries listed above.\n\
         Handle all error cases with AppError variants.\n\
         Do not add new dependencies or imports.",
        function.signature,
        function.available_types,
        function.available_queries,
        function.business_rule_description,
        crystal.implementation.body_patterns.join("\n"),
        crystal.implementation.error_cases.join("\n"),
    )
}
\end{lstlisting}

The prompt gives the LLM:
\begin{enumerate}
\item The exact function signature (constraint: do not change)
\item The types it can use (from Pass 0 models)
\item The queries it can call (from Pass 0 database layer)
\item The business rule description (from TOML, natural language)
\item Known patterns from the crystal (accumulated from past runs)
\item Known error cases (accumulated from past runs)
\end{enumerate}

The LLM operates within a tight box. It fills one function at a time.
It cannot hallucinate types or queries that don't exist because the
prompt lists exactly what is available.

After the LLM returns a function body, it is:
\begin{enumerate}
\item Inserted into the source file (replacing the template body)
\item Checked by isls-reader (Barbara): does it parse?
\item Checked structurally: does it use only declared types/queries?
\item If check fails: retry with error context (up to 3 times)
\item If check passes: move to next function
\end{enumerate}

\subsection{Pass 2: Edge Cases (LLM)}

Takes the code from Pass 1 and asks the LLM to harden it:

\begin{lstlisting}[language=C]
fn build_edge_case_prompt(
    function: &FunctionArtifact,
    crystal: &Crystal,
) -> String {
    format!(
        "Review this Rust function and add edge case handling.\n\
         \n\
         ## Current Implementation\n\
         ```rust\n{}\n```\n\
         \n\
         ## Known Edge Cases for This Pattern\n\
         {}\n\
         \n\
         Consider:\n\
         - What if the input is empty, zero, negative, or very large?\n\
         - What if a referenced entity does not exist?\n\
         - What if a concurrent modification occurs?\n\
         - What if the database connection drops mid-operation?\n\
         - What if the user does not have permission?\n\
         \n\
         Return the COMPLETE function with edge cases added.\n\
         Do not remove existing logic. Only add checks.",
        function.current_body,
        crystal.implementation.edge_cases.join("\n"),
    )
}
\end{lstlisting}

\subsection{Pass 3: Integration (LLM)}

Reviews cross-module consistency:

\begin{lstlisting}[language=C]
fn build_integration_prompt(
    api_file: &str,
    service_file: &str,
    frontend_api_client: &str,
) -> String {
    format!(
        "Check these three files for consistency.\n\
         \n\
         ## API Endpoints (Rust)\n\
         ```rust\n{}\n```\n\
         \n\
         ## Service Functions (Rust)\n\
         ```rust\n{}\n```\n\
         \n\
         ## Frontend API Client (JavaScript)\n\
         ```javascript\n{}\n```\n\
         \n\
         List any mismatches:\n\
         - API route that calls a non-existent service function\n\
         - Frontend URL that doesn't match an API route\n\
         - Request body shape that doesn't match the endpoint's expected payload\n\
         - Response shape that doesn't match what the frontend expects\n\
         \n\
         For each mismatch, provide the fix (which file, which line, what to change).\n\
         Return JSON: [{{\"file\": ..., \"line\": ..., \"fix\": ...}}, ...]",
        api_file, service_file, frontend_api_client,
    )
}
\end{lstlisting}

\subsection{Pass 4: Test Generation (LLM)}

For each service function, generate meaningful tests:

\begin{lstlisting}[language=C]
fn build_test_prompt(
    function: &FunctionArtifact,
    crystal: &Crystal,
) -> String {
    format!(
        "Write Rust integration tests for this function.\n\
         \n\
         ## Function\n\
         ```rust\n{}\n```\n\
         \n\
         ## Test Scenarios to Cover\n\
         {}\n\
         \n\
         ## Additional Scenarios\n\
         - Happy path: normal input, expected output\n\
         - Validation: invalid input, expect error\n\
         - Not found: reference nonexistent entity\n\
         - Authorization: insufficient role, expect 403\n\
         - Idempotency: same operation twice, expected behavior\n\
         \n\
         Use #[tokio::test] and sqlx::PgPool.\n\
         Each test is self-contained: create test data, execute, assert, cleanup.",
        function.current_body,
        crystal.implementation.test_scenarios.join("\n"),
    )
}
\end{lstlisting}

\subsection{Pass 5: Polish (LLM)}

Final pass: improve error messages, add logging, improve docs.
Smallest scope, lowest token budget, tightest convergence threshold.
This pass often converges immediately because Passes 1--4 already
produced clean code.

\subsection{Convergence Measurement}

\begin{lstlisting}[language=C]
/// Measure how much changed between before and after a pass.
/// Returns a value in [0.0, 1.0] where 0.0 = no change.
fn measure_change_rate(
    before: &[Artifact],
    after: &[Artifact],
) -> f64 {
    let mut total_lines = 0usize;
    let mut changed_lines = 0usize;

    for (b, a) in before.iter().zip(after.iter()) {
        if let (Artifact::SourceFile { content: bc, .. },
                Artifact::SourceFile { content: ac, .. }) = (b, a)
        {
            let b_lines: Vec<_> = bc.lines().collect();
            let a_lines: Vec<_> = ac.lines().collect();
            total_lines += a_lines.len().max(b_lines.len());
            // Simple line-by-line diff
            for i in 0..a_lines.len().max(b_lines.len()) {
                let bl = b_lines.get(i).unwrap_or(&"");
                let al = a_lines.get(i).unwrap_or(&"");
                if bl != al { changed_lines += 1; }
            }
        }
    }

    if total_lines == 0 { 0.0 } else { changed_lines as f64 / total_lines as f64 }
}
\end{lstlisting}

% ============================================================
\section{Oracle Integration}

\subsection{Oracle Trait}

\begin{lstlisting}[language=C]
/// The Oracle trait. Abstracts over LLM providers.
pub trait Oracle: Send + Sync {
    /// Send a prompt, get a response.
    fn call(&self, prompt: &str, max_tokens: u32) -> Result<String>;

    /// Send a prompt, get structured JSON response.
    fn call_json(&self, prompt: &str, max_tokens: u32) -> Result<serde_json::Value>;

    /// Name of the model
    fn model_name(&self) -> &str;

    /// Estimated cost per 1K tokens (input + output)
    fn cost_per_1k_tokens(&self) -> f64;
}
\end{lstlisting}

\subsection{OpenAI Implementation}

\begin{lstlisting}[language=C]
pub struct OpenAiOracle {
    api_key: String,
    model: String,      // "gpt-4o-mini" (default), "gpt-4o", etc.
    base_url: String,   // "https://api.openai.com/v1"
    client: reqwest::blocking::Client,
}

impl OpenAiOracle {
    pub fn new(api_key: String) -> Self {
        Self {
            api_key,
            model: "gpt-4o-mini".into(),
            base_url: "https://api.openai.com/v1".into(),
            client: reqwest::blocking::Client::new(),
        }
    }

    pub fn with_model(mut self, model: &str) -> Self {
        self.model = model.into();
        self
    }
}

impl Oracle for OpenAiOracle {
    fn call(&self, prompt: &str, max_tokens: u32) -> Result<String> {
        let body = json!({
            "model": self.model,
            "messages": [
                {"role": "system", "content": "You are a precise code generator. Return only the requested code, no markdown fences, no explanation."},
                {"role": "user", "content": prompt}
            ],
            "max_tokens": max_tokens,
            "temperature": 0.2,  // low temperature for code
        });

        let resp = self.client
            .post(format!("{}/chat/completions", self.base_url))
            .header("Authorization", format!("Bearer {}", self.api_key))
            .header("Content-Type", "application/json")
            .json(&body)
            .send()?;

        let json: serde_json::Value = resp.json()?;
        let content = json["choices"][0]["message"]["content"]
            .as_str()
            .ok_or_else(|| anyhow!("No content in response"))?;

        Ok(content.to_string())
    }

    fn call_json(&self, prompt: &str, max_tokens: u32) -> Result<serde_json::Value> {
        let text = self.call(prompt, max_tokens)?;
        // Strip markdown fences if present
        let clean = text.trim()
            .trim_start_matches("```json")
            .trim_start_matches("```")
            .trim_end_matches("```")
            .trim();
        Ok(serde_json::from_str(clean)?)
    }

    fn model_name(&self) -> &str { &self.model }
    fn cost_per_1k_tokens(&self) -> f64 { 0.15 }  // gpt-4o-mini approximate
}
\end{lstlisting}

\subsection{Mock Oracle (for testing without API key)}

\begin{lstlisting}[language=C]
pub struct MockOracle;

impl Oracle for MockOracle {
    fn call(&self, prompt: &str, _max_tokens: u32) -> Result<String> {
        // Return the existing function body unchanged
        // (effectively a no-op for the render pass)
        if prompt.contains("Return ONLY the function body") {
            Ok("// TODO: implement business logic".into())
        } else if prompt.contains("List any mismatches") {
            Ok("[]".into())
        } else {
            Ok("// mock oracle: no-op".into())
        }
    }

    fn call_json(&self, _prompt: &str, _max_tokens: u32) -> Result<serde_json::Value> {
        Ok(json!([]))
    }

    fn model_name(&self) -> &str { "mock" }
    fn cost_per_1k_tokens(&self) -> f64 { 0.0 }
}
\end{lstlisting}

% ============================================================
\section{CLI Integration}

\begin{lstlisting}[language=bash,title=CLI Commands]
# v2.0 decomposition only (no LLM, same as before)
isls forge-v2 --requirements examples/warehouse.toml \
              --output ./output/warehouse \
              --mock-oracle --trace

# v2.1 full render (with LLM passes)
isls forge-v2 --requirements examples/warehouse.toml \
              --output ./output/warehouse \
              --api-key $OPENAI_API_KEY \
              --model gpt-4o-mini \
              --passes 5 \
              --token-budget 100000 \
              --trace

# v2.1 with specific pass control
isls forge-v2 --requirements examples/warehouse.toml \
              --output ./output/warehouse \
              --api-key $OPENAI_API_KEY \
              --passes 3 \
              --skip-pass polish \
              --trace

# v2.1 render additional passes on existing output
isls render --output ./output/warehouse \
            --api-key $OPENAI_API_KEY \
            --passes 2 \
            --scope services \
            --trace

# Crystal management
isls crystals list
isls crystals stats
isls crystals show belongs_to
isls crystals export --output crystals.json
isls crystals import --file crystals.json
\end{lstlisting}

\subsection{Environment Variable Support}

\begin{lstlisting}[language=bash]
# API key can be set via environment
export OPENAI_API_KEY=sk-...
isls forge-v2 --requirements examples/warehouse.toml --output ./out

# Or via .env file in the working directory
echo "OPENAI_API_KEY=sk-..." > .env
isls forge-v2 --requirements examples/warehouse.toml --output ./out
\end{lstlisting}

% ============================================================
\section{Expected Output with LLM Passes}

\begin{lstlisting}[title=forge-v2 with --passes 5 --trace]
=== ISLS v2.1 Multi-Pass Rendering ===

Pass 0: Structure (offline)
  Files: 60, LOC: 5,120
  Tokens: 0
  Time: 0.8s

Pass 1: Domain Logic (gpt-4o-mini)
  Scope: 12 service functions
  Functions enriched: 12/12
  Tokens: 18,420
  Retries: 2 (parse errors, auto-fixed)
  LOC delta: +1,840
  Time: 34s

Pass 2: Edge Cases (gpt-4o-mini)
  Scope: 12 service functions
  Functions hardened: 9/12 (3 already sufficient)
  Tokens: 8,230
  LOC delta: +620
  Time: 18s

Pass 3: Integration (gpt-4o-mini)
  Scope: API + Services + Frontend
  Mismatches found: 3
  Fixes applied: 3/3
  Tokens: 4,100
  LOC delta: +45
  Time: 8s

Pass 4: Tests (gpt-4o-mini)
  Scope: test files
  Test functions generated: 24
  Tokens: 12,650
  LOC delta: +890
  Time: 22s

Pass 5: Polish (gpt-4o-mini)
  Scope: all
  Improvements: 18 (error messages, logging, comments)
  Tokens: 3,800
  Convergence: 0.008 (< 0.01 threshold -- converged)
  LOC delta: +210
  Time: 7s

=== Summary ===
  Total files: 63
  Total LOC: 8,725
  Total tokens: 47,200
  Estimated cost: $0.007 (gpt-4o-mini)
  Total time: 89s
  Crystals updated: 8

=== Crystal Learning ===
  state_machine: +3 error cases, +2 test scenarios (from order service)
  belongs_to: +1 edge case (cascade delete with active references)
  crud_entity: +2 body patterns (optimistic locking, soft delete check)
\end{lstlisting}

\subsection{Second Run: Different Domain (E-Commerce)}

\begin{lstlisting}[title=ecommerce.toml -- completely different domain]
=== ISLS v2.1 Multi-Pass Rendering ===

Pass 0: Structure (offline)
  Domain detected: ecommerce (keywords: shop, cart, checkout)
  Universal crystals applied: 12/15
  Files: 72, LOC: 6,100
  Tokens: 0

Pass 1: Domain Logic (gpt-4o-mini)
  Functions enriched: 16/16
  Tokens: 22,100
  NOTE: state_machine crystal provided 3 error cases from
        warehouse run -- applied to order status transitions
  NOTE: belongs_to crystal provided cascade edge case --
        applied to cart item deletion

Pass 2: Edge Cases
  Tokens: 5,400 (DOWN from 8,230 -- crystals already had 4 edge cases)

[...]

=== Summary ===
  Total tokens: 41,800 (DOWN from 47,200 on first-ever run)
  Crystal reuse: 12 universal crystals applied
  Cross-domain learning: 7 implementation patterns transferred
\end{lstlisting}

The key metric: \textbf{the second run of a completely different
domain already benefits from universal crystal knowledge gained in
the first run.} Not because the domains are similar. Because the
architectural patterns are universal.

% ============================================================
\section{Implementation Plan}

\subsection{New Crate: isls-renderloop}

\begin{lstlisting}[language=bash]
[dependencies]
isls-types = { path = "../isls-types" }
isls-decomposer = { path = "../isls-decomposer" }
isls-blueprint = { path = "../isls-blueprint" }
isls-orchestrator = { path = "../isls-orchestrator" }
isls-reader = { path = "../isls-reader" }
isls-code-topo = { path = "../isls-code-topo" }
isls-evidence = { path = "../isls-evidence" }
serde = { version = "1", features = ["derive"] }
serde_json = "1"
reqwest = { version = "0.12", features = ["json", "blocking"] }
sha2 = "0.10"
anyhow = "1"
tracing = "0.1"
dotenv = "0.15"
\end{lstlisting}

Contains: RenderLoop, RenderPass, PassType, Oracle trait,
OpenAiOracle, MockOracle, convergence measurement, prompt
builders for each pass type.

\subsection{Modified: isls-blueprint}

Restructure Crystal type to include CrystalLevel and
ImplementationKnowledge. Add 15 built-in universal crystals.
Change matching logic to search by level (universal first).

\subsection{Modified: isls-decomposer}

After decomposition, hand artifacts to RenderLoop instead of
returning directly. The decomposer becomes Pass 0.

\subsection{Modified: isls-cli}

Add: \texttt{--api-key}, \texttt{--model}, \texttt{--passes},
\texttt{--token-budget}, \texttt{--skip-pass}, \texttt{--scope}.
Add: \texttt{isls render} subcommand for re-rendering existing output.
Add: \texttt{isls crystals} subcommand group.

\subsection{Execution Order}

\begin{enumerate}
\item Restructure \texttt{isls-blueprint}: CrystalLevel,
  ImplementationKnowledge, 15 universal crystals, level-based
  matching. Test in isolation.

\item Create \texttt{isls-renderloop}: Oracle trait, OpenAiOracle,
  MockOracle. Test: mock oracle call returns string.

\item Implement Pass 1 (DomainLogic): prompt builder, function
  body insertion, Barbara parse-check. Test: mock oracle enriches
  one function, file still parses.

\item Implement Pass 2 (EdgeCases): prompt builder, edge case
  insertion. Test: mock adds edge case comment to function.

\item Implement Pass 3 (Integration): cross-module check.
  Test: detect deliberate mismatch in test fixture.

\item Implement Pass 4 (TestGeneration): test prompt builder.
  Test: mock generates test function skeleton.

\item Implement convergence measurement. Test: identical
  before/after = 0.0 change rate.

\item Wire RenderLoop into isls-decomposer pipeline.
  \texttt{forge-v2 --mock-oracle} still works (Pass 0 only).
  \texttt{forge-v2 --api-key X} runs all passes.

\item Implement crystal learning: after Pass 1--4, extract
  implementation knowledge, update crystals. Test: crystal
  gains body\_patterns after mock run.

\item CLI: all new flags and subcommands.

\item End-to-end test with mock oracle: full render loop,
  all passes execute, convergence reached, crystals updated.

\item (Optional) End-to-end test with real API key if available.
\end{enumerate}

% ============================================================
\section{Success Criteria}

\begin{enumerate}
\item \texttt{cargo build --workspace} --- zero errors.
\item All existing tests pass (v2.0 unchanged).
\item \texttt{forge-v2 --mock-oracle} still works (Pass 0 only).
\item \texttt{forge-v2 --mock-oracle --passes 5} runs all 6 passes
  with mock oracle, produces output, reports convergence.
\item With real API key: \texttt{forge-v2 --api-key \$KEY --passes 5}
  enriches service functions with actual business logic.
\item Crystal registry contains 15 built-in universal crystals.
\item Crystal matching searches universal $\to$ architectural $\to$
  structural $\to$ domain (in that order).
\item After a render run, crystals have updated ImplementationKnowledge
  (body\_patterns, error\_cases, edge\_cases).
\item Convergence measurement works: Pass 5 shows lower change rate
  than Pass 1.
\item Token usage is tracked per pass and reported in trace output.
\item Second run with different domain TOML benefits from universal
  crystal knowledge (reported in trace).
\end{enumerate}

% ============================================================
\section{Implementation Constraints for Claude Code}

\begin{enumerate}
\item \textbf{Repository: graphity/isls.} Add isls-renderloop
  to workspace members.

\item \textbf{v2.0 must not break.} \texttt{forge-v2 --mock-oracle}
  without \texttt{--passes} or \texttt{--api-key} must produce
  exactly the same output as before.

\item \textbf{reqwest blocking.} Use \texttt{reqwest::blocking}
  for Oracle HTTP calls. No async runtime needed --- the render
  loop is sequential (one prompt at a time).

\item \textbf{API key from environment or CLI.} Support both
  \texttt{--api-key} flag and \texttt{OPENAI\_API\_KEY} env var.
  Use \texttt{dotenv} crate to load \texttt{.env} file.

\item \textbf{Token counting.} Estimate tokens as
  \texttt{text.len() / 4} (rough approximation). Do not add
  a tokenizer dependency.

\item \textbf{Prompt construction is critical.} The prompts in
  this spec are the exact templates. Do not simplify them. The
  quality of LLM output depends entirely on prompt quality.

\item \textbf{All 15 universal crystals must be defined with
  complete StructuralKnowledge.} The ImplementationKnowledge
  starts empty --- it is filled by the render loop.

\item \textbf{Read the existing isls-blueprint code} before
  restructuring. Preserve backward compatibility.

\item \textbf{Sequential implementation.} Blueprint restructure
  $\to$ renderloop crate $\to$ passes $\to$ convergence $\to$
  CLI $\to$ end-to-end.

\item \textbf{The mock-oracle end-to-end test is the most
  important deliverable.} All 6 passes run, convergence is
  reported, crystals are updated. If time runs out, skip
  the real-API test but keep mock working.

\item \textbf{No \texttt{unwrap()} in library code.}
\item \textbf{Every public item documented.}
\item \textbf{MIT License. Copyright (c) 2026 Sebastian Klemm.}
\end{enumerate}

\vfill
\noindent\rule{\textwidth}{0.4pt}
\begin{center}
\small
ISLS v2.1 --- Multi-Pass Rendering --- Sebastian Klemm --- March 2026\\
Structure offline. Logic by LLM. Edge cases by LLM. Converge. Learn.\\
Universal crystals learn from every domain. Every run makes the next better.\\
Not for one domain. For all of them.\\
\url{https://github.com/LashSesh/graphity}
\end{center}

\end{document}
